\chapter{Verwandte Ideen und mögliche Folgen}
\section{Zuse und Wolfram: eine Verwandtschaft ohne Rangordnung}
Der eingefügte Gesprächstext vergleicht den Universalraum mit Konrad Zuse und Stephen Wolfram. Zuses Werk heißt \emph{Rechnender Raum}, nicht „denkender Raum“. In seinem Buch von 1969 untersucht er digitale Beschreibungen physikalischer Felder und Teilchen und die Rolle von Automaten. Das ist ein historischer Bezugspunkt für die Vorstellung von Physik als regelgeleitetem Prozess. \href{https://link.springer.com/book/10.1007/978-3-663-02723-2}{Zuse, Rechnender Raum [E17]}

Wolframs Programm arbeitet mit der lokalen Umschreibung von Hypergraphen. Auch dort sollen räumliche Kontinuität und Dimension erst als große Muster entstehen; Zeit wird mit Update-Ereignissen verbunden. Multiway-Strukturen und kausale Invarianz sollen verschiedene mögliche Updatefolgen organisieren. Das sind Ziele und vorgeschlagene Mechanismen seines Programms, keine hier bestätigten universellen Ableitungssätze. \href{https://www.wolframphysics.org/technical-introduction/potential-relation-to-physics/basic-concepts/}{Wolfram Physics, Basic Concepts [E18]}

Die Formulierung im Gesprächstext, TFPT gehe automatisch „eine Ebene tiefer“, ist daher keine belastbare wissenschaftliche Abgrenzung. Auch Graph-Umschreibungen besitzen eine Algebra ihrer Kompositionen, und auch Wolfram postuliert nicht einfach einen vorher gegebenen gewöhnlichen Raum. Der konkrete Unterschied liegt in den gewählten Ausgangsdaten und Nachweisen:
\par\noindent\begin{tabularx}{\linewidth}{L{26mm}Y Y}
\toprule \textbf{Ansatz}&\textbf{Betonte Ausgangsstruktur}&\textbf{Für einen Vergleich zu prüfen}\\\midrule
Zuse&Digitale Vorgänge und Automaten&Welche diskrete Regel erzeugt welche physikalische Entwicklung?\\
Wolfram&Hypergraph-Umschreibungen und Updategeschichten&Wie werden Kausalstruktur, Dimension und Quanteneigenschaften gewonnen?\\
TFPT/Universalraum&Markierte Algebra, zulässige Operationen, positive Prozessstatistik&Wie wählt dieselbe Quelle Kopplung, Zustand, Register und physikalischen Grenzprozess?\\\bottomrule
\end{tabularx}\par
Die Gemeinsamkeit ist die Suche nach Struktur aus Vorgängen. Die Modelle werden erst durch explizite Regeln, kontrollierte Grenzwerte und unterscheidende Vorhersagen vergleichbar.

\section{Was schon heute praktisch nutzbar ist}
Die unmittelbaren Ergebnisse sind Werkzeuge zur Forschung: genaue Tests für verlorene Gedächtnisinformation; kleine Quantenmodelle, deren Energie und Aufzeichnung gemeinsam berechenbar sind; Zustandszertifikate; Gegenbeispiele gegen falsche Eindeutigkeitsannahmen; und ein Schaltungsprotokoll mit vollständiger Erfolgsrechnung.

Methodisch lässt sich daraus auch außerhalb der Fundamentalphysik lernen: Gleiche Antworten jetzt bedeuten nicht gleiches Verhalten später. Man muss den Zugriff auf gespeicherte Information und die erlaubten Eingriffe prüfen. Die im Korpus erwähnte Übertragung auf Hylæan ist genau diese methodische Analogie, keine aus dem Universalraum abgeleitete Fähigkeit eines KI-Systems.

\section{Was eine vollständige Theorie verändern könnte}
Falls ein gemeinsamer Ursprung tatsächlich rekonstruiert würde, könnte er bisher getrennte Parameterbeziehungen erklären, neue konsistente Materiesektoren eingrenzen und präzisere Aussagen über Grenzen von Simulation und Steuerung erlauben. Daraus könnten irgendwann konkrete Material-, Mess- oder Rechenverfahren entstehen. Jede solche Anwendung benötigt zusätzlich einen zugänglichen Betriebsbereich, einen steuerbaren Effekt und eine Kostenrechnung.

Die im Gesprächstext erwähnten Ideen künstlicher Gravitation, lokaler Geometrieänderung, anderer Vakuumzustände oder ungewöhnlicher Raumzeitwege sind viel weitergehende Spekulationen. Aus den Anhängen folgt kein steuerbarer Mechanismus dafür. Eine vollständige Theorie könnte solche Möglichkeiten auch ausschließen. Verstehen und technisch herstellen sind verschiedene Leistungen; eine emergente Geometrie ist nicht automatisch beliebig verformbar.

\section{RH, Faktorisierung und P versus NP}
Die Universalraumlinie liefert keinen neuen Beweis dieser Probleme. Die in S04 genannten RH-Katalogzahlen -- 2753 Records, 1097 kuratierte, 122 offene, 319 kuratierte beziehungsweise 405 gesamte Kills -- sind ein datierter Bericht dieser Quelle und hier nicht live geprüft.

Die wichtigen negativen Befunde sind inhaltlich:
\begin{itemize}
\item \textbf{RH:} $E_8$-Kongruenzdaten oder eine Zählfunktion aus verschobenen $\zeta$-Funktionen liefern keinen selbstadjungierten Operator mit dem benötigten Nullstellenspektrum. $\prod_{k=0}^7\zeta(s-k)$ enthält die ursprüngliche Nullstellenfrage, löst sie aber nicht. Ein Hecke-/Jones-Index ist noch keine Hilbert-Pólya-Spektralbrücke.
\item \textbf{Faktorisierung:} S04 berichtet die Gaußsummenidentität $S_N(t)=N^4\gcd(t,N)^4$ für den dortigen Ansatz. Die Faktoren können in einer Zählung kodiert sein, während ihre Auslesung teuer bleibt. Die berichtete $O(N^3)$-Route ist in der Eingabelänge $\log N$ nicht polynomial.
\item \textbf{Rechenkomplexität:} Reine Clifford-Prozesse mit Stabilizer-Eingängen und den entsprechenden Messungen sind klassisch effizient simulierbar. Diese Aussage darf nicht pauschal auf $\Omega$-Präparation, allgemeine kontinuierliche Austauschdynamik oder zusätzliche Nicht-Clifford-Ressourcen ausgedehnt werden. Deren Existenz beweist aber umgekehrt weder skalierbare Rechenuniversalität noch einen Faktorisierungsalgorithmus oder P=NP.
\end{itemize}
Die zugeordneten Katalog-IDs in S04 sind r618, r613, r604, r647, TFPT.HECKE.INDEX.01, PRIME.HECKEMODRAM.01 und E8.COXETER.EULER.COMPLETION.01. Ihre Dossiers wurden hier nicht erneut auditiert. Lebende RH-Fronten, die S04 getrennt nennt, sind WEIL\_POSITIVITY\_WINDOWS und SCREW\_SUBORDINATION\_LSTAR; sie sind keine durch dieses Buch geschlossenen Universalraum-Routen.

\chapter{Was noch offen ist und wie es weitergehen sollte}
\label{ch:offen}
\begin{bild}{Aktueller Stand v1.3}
Der endliche Kern ist jetzt konkreter: Ein gemeinsamer korrigierter Vermittler liefert Austauschregister, kontrollierte Präparation und Mehrzeitantwort. Die entscheidende offene Stelle ist seine native nicht-Cliffordsche Realisierung einschließlich kohärentem Belegungszugriff. Präparation per Messung und Reset ist weiterhin vom autonomen physikalischen Zustand verschieden. Die folgende Gesamtkarte gilt mit dieser Ergänzung; die neuen Abnahmetests stehen in Kapitel~\ref{ch:execution13}.
\end{bild}
\section{Die fünf Fugen nach allen Korrekturen}
\begin{longtable}{L{31mm}L{62mm}L{57mm}}
\toprule \textbf{Fuge}&\textbf{Erreichter Stand}&\textbf{Verbleibende Herkunft}\\\midrule\endhead
Kopplungsstärke&Exakter antisymmetrischer Vermittler; $J_{\mathrm{eff}}>0$ unter benannten Bedingungen; S11 liefert $2t^2/\Delta$ in seiner Normierung.&Energien, Amplituden, konkurrierende Kanäle und ihre gemeinsame Auswahl.\\
Kopplungsgraph&Clebsch-Graph exakt unter „Ort = Spinorgewicht“; Tetramer, Kette und Clebsch getrennt.&Ortslesart, Belegung, Lokalität und skalierende Familie von Graphen.\\
Cross-Register-Kopplung&Vormessung und kontrollierter Swap ausdrücklich schreibbar; pauschaler Sektor-Unmöglichkeitsschluss korrigiert.&Warum die Quelle trägerübergreifende Kontrolle und genau diese Ausführung bereitstellt.\\
Präparation&$\Omega$-Schaltung mit Erfolgsanzeige und quantifiziertem Aufwand; Austausch-Erhaltungsgröße sichtbar.&Nicht-Stabilizer-Ressource, Eingangs- und Hilfszustände sowie deren physischer Ursprung.\\
Zeitskala&Mehrere lesbare Uhrprotokolle und ein Signal derselben Zweizellendynamik.&Gemeinsame Energie-/Zeitreferenz, Auswahl des Uhrwerks und Anfangskorrelationen.\\\bottomrule
\end{longtable}
Hinzu kommen Kontextgewichte, das positive Quellfunktional und die physikalische Auslesung. Die fünf Fugen sind also eine nützliche lokale Arbeitskarte, nicht der vollständige TOE-Vertrag.

\section{T1 bis T8: der vollständige Vertrag}
\begin{longtable}{L{9mm}L{62mm}L{79mm}}
\toprule \textbf{Tor}&\textbf{Was gemeinsam nachzuweisen ist}&\textbf{Was bisher dazu vorliegt}\\\midrule\endhead
T1&Herkunft von P1/P2, Dimensionen, Markierungen und Auswahlregeln.&Exakte Konsequenzen festgelegter Daten; Gegenmodelle gegen Symmetrie-Eindeutigkeit. Herkunft offen.\\
T2&Native half-charge-markierte $E_8$-Naht samt Skalierungslimes, Domänen, Energieabschätzungen und beiden Adjungierten.&Gitterhülle und Zielerweiterung konkret; endliche Kontexte und Sektorprüfungen. Native analytische Schließung offen.\\
T3&Ein ausgewählter lokaler oder quasilokaler unitärer Ursprung in 3+1 Dimensionen.&Endliche Kopplungen und bedingte Graphlesarten. Räumlicher Ursprung und gemeinsamer Kegel offen.\\
T4&Chirales Standardmodellmaß, Anomalien/Index, Spiegelentkopplung.&Ladungstabelle und endliche Anomaliesummen. Propagierendes chirales Maß offen.\\
T5&Wechselwirkender Kontinuumsgrenzwert mit Streuung und Clusterstruktur.&Endliche Spektren, bekannte Kettenanschlüsse, genaue kleine Sektoren. Gemeinsamer physischer Grenzwert offen.\\
T6&Alle Eichkopplungen und vollständige Neutrinotextur/-skala aus dem Inneren.&Präzise einzelne Formeln und bedingte Transfers. Vollständige Auswahl offen.\\
T7&Quantisierter masseloser Spin 2, zwei Helizitäten, universelle Kopplung.&Lorentzkinematik und bedingte Zielarchitekturen. Dynamischer gravitativer Sektor offen.\\
T8&Physikalischer Zustand und ein gemeinsames Quellfunktional aller Ausgaben.&$\Omega$, bedingte Präparation, GNS-/Prozessrahmen und Registerzeugen. Ursprüngliche Zustandswahl offen.\\\bottomrule
\end{longtable}
Keines dieser acht Tore wird durch die vorliegenden Anhänge oder die neuen endlichen Gegenrechnungen vollständig geschlossen. Das sagt nichts gegen die Gültigkeit der genauen Teilergebnisse; es beschreibt die größere Reichweite des geforderten Satzes.

\section{Ein priorisiertes, entscheidbares Anschlussprogramm}
\textbf{1. Eine einzige primitive Quelle festlegen.} Sie muss Träger, Vertizes, Phasen, Adjunktion, Belegung, Kopplungszugriffe, Zustandsbedingungen und zulässige Auslesungen enthalten. Jede nicht hergeleitete Eingabe bleibt sichtbar. Erfolg ist ein ausführbarer Quellenvertrag; Scheitern ist eine noch frei wählbare Regel, die das Ergebnis verändert.

\textbf{2. Den Kopplungsmechanismus im selben Modell prüfen.} Die korrekten $\Z_4$-Typen müssen eingehalten, konkurrierende Vermittler und Mehrkörperterme kontrolliert werden. Ein Clebsch-Ursprung muss seine Belegungsstruktur erklären. Der Vergleich mit der Zweizellenformel ist dann eine präzise endliche Kontrolle, falls der Kandidat gerade dieses Limit besitzt; ein anderer nativer Graph darf nicht künstlich auf diese Formel zugeschnitten werden.

\textbf{3. Präparation und Mehrzeitauslesung aus derselben Quelle gewinnen.} Die Quelle muss die Nicht-Stabilizer-Ressource und das Registerprotokoll bestimmen. Der originale Vergleich $1$ gegen $17/32$ und die neue ausgeglichene Variante $1$ gegen $1/2$ sind konkrete mögliche Zeugen. Erfolg ist eine vorher festgelegte Rohwahrscheinlichkeit samt allen Misserfolgen, die ohne nachträgliche Registerwahl berechnet wird.

\textbf{4. Die chirale Naht herstellen.} Nicht nur $c=8$, sondern Gewicht-1-Inhalt, $\Z_4$-Sektoren, Operatorprodukte, Phasen und linke/rechte Struktur müssen gemeinsam passen. Die 248 Zustände sind ein früher Test, kein vollständiger Abschluss. Die Quelle muss auch die zusätzlichen $D_5$-Moden und ihre Randbedingungen tragen.

\textbf{5. Geometrie aus tatsächlicher Ausbreitung ableiten.} Man berechnet die Korrelationsfunktionen desselben Prozesses, ihre Pole, Volumen- und Spektralskalierung sowie die Kegel verschiedener Materiesektoren. Eine bereits vorgegebene Weylmatrix darf nicht anschließend als abgeleitete Raumzeit ausgegeben werden. Unterschiedliche unvereinbare Kegel oder fehlende kohärente Pole wären informative negative Ergebnisse.

\textbf{6. Eine unabhängige physikalische Ausgabe einfrieren.} Erst Inputs, Schema, Transfer, Unsicherheit und Ablehnungskriterium festlegen; dann mit Daten vergleichen. Die ungünstigen Lepton-, Higgs-, Inflations- und Protonzweige bleiben erhalten. Eine externe Planckskala zählt nicht nochmals als unabhängige Vorhersage.

\section{Die präziseste gemeinsame Hypothese}
\begin{bild}{Der Kandidat, auf den die Dokumente gemeinsam zulaufen}
Es gibt eine phasentreue Kompositionsregel mit positivem Zustandsfunktional, deren endliche Auslesungen die markierte TFPT-Struktur tragen. Aus ihren gekoppelten Anregungen und Aufzeichnungen entsteht ein kontrollierter gemeinsamer Grenzprozess. Erst dieser Grenzprozess liefert Raumzeit, chirale Materie, Wechselwirkungen und gegebenenfalls Gravitation.
\end{bild}
Die erste Hälfte ist an mehreren kleinen Stellen konkret realisiert. Die Auswahl der gemeinsamen Quelle und der vollständige physikalische Grenzprozess sind noch offen. Die neue Synthese verkleinert die Lücke dort, wo sie einen falschen Eindeutigkeitsschluss ersetzt, eine Kopplung ausdrücklich konstruiert oder eine neue unterscheidende Messung berechnet. Genau daran lässt sich der nächste Fortschritt messen.

\par\noindent\textbf{Fortsetzung in Version 1.1:} Die folgenden fünf Kapitel untersuchen zusätzliche Schatten-Phänomene. Die zusammengeführte Arbeitshypothese und der neue endliche Interferenztest stehen in Kapitel~\ref{ch:shadow-synthesis}.
